\section{Introduction}
\label{sec:introduction}

Troubleshooting computer networks is a time-consuming and still largely manual
activity. When a service reports a problem or a host loses connectivity, an
engineer needs to consult the available tools and check one device after another, running
show commands, inspecting interface counters, routing tables, ARP caches
and syslog messages until a fault is localised. The path from the first symptom to the true root cause
often requires a lot of steps and time. As networks
grow in size and complexity, this diagnostic effort translates directly into
operational cost and into mean-time-to-repair.

Recent advances in large language models (LLMs) and, in particular, in
agentic systems that let an LLM call external tools, observe the results
and decide on the next action, suggest a new way to support this work. Instead of
a single agent model, a multi-agent system can decompose the problem
into smaller parts and assign them to specialised agents. 
Whether such a system can be used to assist in troubleshooting or even find problems autonomously is an open
and practically relevant question.

\subsection{Project goal}

The goal of this project is to design and build an AI-driven multi-agent system
for network troubleshooting and to use it as an instrument to evaluate
\emph{which kinds of network problems can realistically be troubleshooted by such
a system, and at what benefit and cost}. 

\subsection{Research questions}
\label{sec:research-questions}

The project is guided by four research questions:

\begin{enumerate}
  \item \textbf{Time}: How much time can be saved in troubleshooting by using the
        agent network compared to a manual workflow?
  \item \textbf{Difficulty}: Problems of which difficulty level can be
        troubleshooted by the agent network?
  \item \textbf{Success rate}: What is the success rate for the correct analysis
        of the problems at each difficulty level?
  \item \textbf{Cost}: How much can potentially be saved or, conversely, how
        much does it cost to use an agent network for troubleshooting network
        issues, and how does this depend on the chosen model?
\end{enumerate}


\subsection{Outline}

\Cref{sec:background} introduces the underlying technologies. LLMs, AI agents,
function calling, the agent loop and the Model Context Protocol.
\Cref{sec:concepts} explains how these building blocks are combined into the
multi-agent troubleshooting system and \cref{sec:implementation} describes its
realisation. \Cref{sec:methodology} defines the fault scenarios, the
difficulty tiers, the evaluated models and the experiment procedure.
\Cref{sec:results} reports the measurements, \cref{sec:evaluation} interprets
them with respect to the research questions and discusses the limitations, and
\cref{sec:conclusion} concludes the report and outlines future work.